\section{Nástroje pro správu verzí softwaru}
Následující podkapitola se věnuje nástrojům pro správu verzí softwaru pomocí verzovacího systému Git. Tyto nástroje však v dnešní době nepodporují pouze samotnou správu verzí softwaru, ale nabízí i ticketovací systémy, možnost tvorby stránek s dokumentací, nástroje pro CI/CD a mnoho dalších souvisejících funkcionalit. V dalších sekcích se budu věnovat pouze nástrojům GitHub a GitLab, ale mezi další zástupce patří například nástroj Bitbucket od firmy Atlassian.

\subsection{GitHub}
Vývoj platformy GitHub začal v roce 2007 a od té doby její uživatelská základna neustále roste. Na začátku roku 2023 už GitHub využívalo více než 100 milionů uživatelů a 4 milióny organizací na správu softwarových projektů. \cite{GitHub}\cite{GitHub-about} Základní funcionality jsou dostupné zdarma jak u veřejných, tak u privátních repozitářů, ale pokročilejší funkce, které souvisejí například se zabezpečením jako jsou automatické detekce úniků tokenů a hesel nebo nastavování pravidel pro zahrnování změn do chráněných větví, jsou dostupné pouze pro veřejné repozitáře nebo jsou součástí placených verzí. Self-hosted řešení je dostupné pouze v placené enterprise variantě. \cite{GitHub-pricing}

\subsection{GitLab}
GitLab začal být vyvíjen jako open source projekt v roce 2011 a od té doby se na jeho vylepšování podílelo více než 3300 přispěvatelů. Aktuálně má GitLab dle odhadů přes 30 milionů registrovaných uživatelů, kteří si mohou stejně jako u GitHubu vybrat z několika nabízených cenových plánů. \cite{GitLab-about} Díky komunitní open source verzi ale GitLab nabízí i bezplatnou self-hosted variantu, která organizacím umožňuje mít všechna data pod svou kontrolou. \cite{GitLab-pricing}

\subsection{APIs}
GitHub i GitLab poskytují REST a GraphQL APIs, pomocí nichž lze provádět nejrůznější integrace, vytvářet nástroje na automatizaci procesů nebo mohou sloužit jako zdroj dat pro analytické dashboardy. Dokumentace těchto APIs jsou dostupné online na webových stránkách GitHubu \cite{GitHub-REST}\cite{GitHub-GraphQL} a GitLabu \cite{GitLab-REST}\cite{GitLab-GraphQL}. Ačkoliv je u veřejných repozitářů možné provádět některé akce bez autorizace, tak je s ohledem na nižší limity pro neautorizované uživatele a další omezení, které mají zabránit nedostupnosti služeb, většinou nutné zasílat informaci o autorizaci při každém požadavku. Dostupné způsoby autorizace se mezi službami GitHub a GitLab a mezi jednotlivými druhy API liší a budou popsány v podkapitole xx.\todo{přidat odkaz}

\begin{description}
    \item[REST] TODO
    \item[GraphQL] TODO
\end{description}

Na níže uvedeném příkladu je vidět, jakým způsobem lze pomocí REST a GraphQL API u nástroje GitHub získat popis repozitáře \mintinline{text}|kuoris-v2| a název prvních deseti issues. V případě REST API by bylo potřeba provést dva GET požadavky -- první na získání popisu repozitáře a druhý pro získání názvů issues viz výpis kódu \ref{listing:rest-query-example}. Oproti tomu by v případě GraphQL API stačilo provést jeden POST podažavek na \mintinline{text}|https://api.github.com/graphql| s tělem, který je vidět na výpisu kódu \ref{listing:graphql-query-example}.
 
\begin{listing}[h]
    \caption{Ukázka získání informací o repozitáři a issues pomocí REST API}\label{listing:rest-query-example}
    \begin{minted}{bash}
curl https://api.github.com/repos/paukert/kuoris-v2
curl https://api.github.com/repos/paukert/kuoris-v2/issues
    \end{minted}
\end{listing}

\begin{listing}[h]
    \caption{Ukázka těla požadavku pro GraphQL API}\label{listing:graphql-query-example}
    \begin{minted}{graphql}
query {
    repository(owner: "paukert", name: "kuoris-v2") {
        description
        issues(first: 10) {
            nodes {
                title
            }
        }
    }
}
    \end{minted}
\end{listing}

Ještě doprozkoumat, ale GraphQL se jeví jako zajímavější možnost.

Popsat možnosti autorizace pro vybraný typ API (asi tedy GraphQL).

\begin{description}
    \item[GitHub] Authenticating with a personal access token
    \item[GitLab]
\end{description}
